% Part: computability
% Chapter: machines-computations
% Section: halting-problem

\documentclass[../../../include/open-logic-section]{subfiles}

\begin{document}

\olfileid{tur}{und}{hal}
\olsection{停机问题}

\begin{explain}
假设我们已经固定了图灵机描述的一种枚举。
因此，每台图灵机都获得一个\emph{指标}：即它在图灵机描述枚举 $M_1$, $M_2$, $M_3$, \dots{} 中的位置。

我们知道，一定存在非图灵可计算的函数：图灵机描述的集合——因而图灵机的集合——是可数的，但从 $\Nat$ 到 $\Nat$ 的所有函数的集合却不可数。不过，我们也能找到不可计算函数的具体例子。一个这样的例子就是停机函数。
\end{explain}

\begin{defn}[停机函数]
 \emph{停机函数}~$h$ 定义为
\[
h(e,n) =
\begin{cases}
  \text{0} & \text{若机器~$M_e$ 在输入 $n$ 上不停机} \\
  \text{1} & \text{若机器~$M_e$ 在输入 $n$ 上停机}
\end{cases}
\]
\end{defn}

\begin{defn}[停机问题]
\emph{停机问题}是指（对任意 $e$, $n$）判定图灵机~$M_e$ 是否在输入为~$n$ 个竖杠时停机的问题。
\end{defn}

\begin{explain}
我们通过证明一个相关的函数~$s$ 不是图灵可计算的，来证明~$h$ 不是图灵可计算的。这个证明依赖于以下事实：任何能被图灵机计算的东西，也都能被一台规范的图灵机计算（\olref[mac][dis]{sec}），以及两台图灵机可以连接起来构成一台单独的机器（\olref[mac][cmb]{sec}）。
\end{explain}

\begin{defn} 函数~$s$ 定义为
\[
s(e) =
\begin{cases}
  \text{0} & \text{若机器~$M_e$ 在输入 $e$ 上不停机} \\
  \text{1} & \text{若机器~$M_e$ 在输入 $e$ 上停机}
\end{cases}
\]
\end{defn}

\begin{lem}
函数~$s$ 不是图灵可计算的。
\end{lem}

\begin{proof}
为了推出矛盾，假设函数~$s$ 是图灵可计算的。那么，就会有一台计算~$s$ 的图灵机~$S$。我们可以不失一般性地假定，当 $S$ 停机时，它停在第一个方格上（即它是规范的）。这台机器可以“连接”到另一台机器~$J$ 上，$J$ 的行为是：如果它在输入为~0 时启动（即，在初始状态下扫描到磁带末端符号右侧的方格为 $\TMblank$），那么它就停机；否则，它就会向右无限移动，永不停机。将 $S$ 连接到~$J$ 所得到的机器 $S \concat J$ 是一台图灵机，因此它是某个~$e$ 对应的 $M_e$（即它出现在枚举中的某个位置）。让 $M_e$ 以 $e$ 个 $\TMstroke$ 为输入启动。有两种可能性：要么 $M_e$ 停机，要么不停机。
\begin{enumerate}
\item 假设 $M_e$ 在输入为 $e$ 个 $\TMstroke$ 时停机。那么 $s(e) = 1$。因此，当 $S$ 以~$e$ 为输入启动时，它会停机，并在磁带上留下单个 $\TMstroke$ 作为输出。接着，当 $J$ 启动时，磁带上有一个 $\TMstroke$。在这种情况下，$J$ 不会停机。但是 $M_e$ 就是机器 $S \concat J$，所以它所做的应该恰好是 $S$ 后接 $J$ 会做的事（即，在这种情形下，向右无限移动且永不停机）。因此，$M_e$ 不可能在输入为 $e$ 个 $\TMstroke$ 时停机。

\item 现在假设 $M_e$ 在输入为 $e$ 个 $\TMstroke$ 时不停机。那么 $s(e) = 0$，并且 $S$ 在以~$e$ 为输入启动时，会停机在空白磁带上。当 $J$ 在空白磁带上启动时，会立即停机。同样，$M_e$ 所做的是 $S$ 后接~$J$ 会做的事，因此 $M_e$ 必须在输入为 $e$ 个 $\TMstroke$ 时停机。
\end{enumerate}
在每种情形下，我们都得到了与假设相矛盾的结论。这表明不可能存在图灵机~$S$：$s$~不是图灵可计算的。
\end{proof}

\begin{thm}[停机问题的不可解性]
\ollabel{thm:halting-problem} 停机问题是不可解的，即函数~$h$ 不是图灵可计算的。
\end{thm}

\begin{proof}
假设 $h$ 是图灵可计算的，即由一台图灵机~$H$ 计算。我们可以用 $H$ 来构造一台计算~$s$ 的图灵机：首先，复制输入（用一个 $\TMblank$ 符号分隔）。然后移回开头，并运行~$H$。我们显然可以构造一台完成前一项操作的机器（参见 \cref{tur:mac:dis:prob:copier}），并且如果 $H$ 存在，我们就能把它“连接”到这样一台复制机上，从而得到一台能判定 $M_e$ 在输入为~$e$ 时是否停机的新机器，也就是说，它能计算~$s$。但我们已经证明了这样的机器不可能存在。因此，$h$~也不是图灵可计算的。
\end{proof}

\begin{prob}
三停机（3-Halt）问题是指，给出一个判定程序，用以判定任意选定的图灵机当输入为三个 $\TMstroke$（其余磁带为空白）时是否停机。证明三停机问题是不可解的。
\end{prob}

\begin{prob}
证明：若对于图灵机与输入的配对 $M_e$ 和~$n$（其中 $e \neq n$），停机问题是可解的，则对于 $e = n$ 的情况，它也是可解的。
\end{prob}

\begin{prob}
我们证明了，如果输入是一个数字~$e$（它通过枚举所有图灵机来标识一台图灵机~$M_e$），那么停机问题是不可解的。如果我们允许直接使用 \olref[tur][und][enu]{sec} 节中对图灵机的描述作为输入，情况会怎样？能否存在一台图灵机，它能够解决停机问题，但以图灵机的描述而非指标作为输入？请说明理由。
\end{prob}

\begin{prob} 证明如下定义的\emph{部分}函数~$s'$ 是图灵可计算的：
  \[
  s'(e) =
  \begin{cases}
    \text{1} & \text{若机器~$M_e$ 在输入 $e$ 上停机}\\
    \text{无定义} & \text{若机器~$M_e$ 在输入 $e$ 上不停机}
  \end{cases}
  \]
\end{prob}

\end{document}